\documentclass{situationist-preprint}

\title{The Recognized Self}
\msid{PID-12-a}
\msfamily{PID-12}
\msdiscipline{Personal Identity and the Self}
\mstype{Principal manuscript}
\mskeywords{self; recognition; social self; practical identity; Hegel; Mead; Honneth; constitution}

\begin{document}
\maketitle

\begin{abstract}
Is the self social? Individualists hold that a self is what it is independently of others, who at most influence it causally; social theorists from Hegel and Mead to Honneth hold that selves are made in relations of recognition. We argue for a qualified answer: the self is partly constituted by relations of recognition. The self, as PID-04 argues, is a construction, a structure of self-representation, self-interpretation and self-regulation. The minimal subject of experience does not depend constitutively on recognition. But the practical identities by which a person interprets and governs herself, as a parent, a member of a profession, a citizen or a reliable friend, are statuses whose holding consists in part in their recognition by others, and the normative standing to have one's self-interpretations count, which narrative identity requires (PID-09), is conferred in the same way. Since these components are parts of the self, the self is partly constituted by recognition; since the others are not, only partly. We answer the objections that recognition is merely causal, that the view makes the self hostage to others, and that it cannot account for the solitary.
\end{abstract}

\section{The problem}

\citet[§§178--196]{Hegel1807} argued that self-consciousness achieves its satisfaction only in another self-consciousness, and that the struggle for recognition is part of the formation of the self. \citet{Mead1934} held that the self arises in social experience, through taking the attitudes of others toward oneself. \citet{Taylor1994} argued that identity is formed in dialogue and that misrecognition can inflict real harm, and \citet{Honneth1995} distinguished three forms of recognition, love, rights and esteem, each corresponding to a form of practical relation to oneself. On the other side, the Cartesian and Lockean traditions take the self to be given to itself in consciousness, prior to any relation to others, and contemporary defenders of a minimal self locate its core in the first-person character of experience \citep{GStrawson1997}.

Each side has a strong case for part of what is called the self. We argue that the self has parts of both kinds, and that this is the sense in which it is social: partly.

\section{Framework}

\begin{definition}[Constitution and causation]\label{d:const}
A relation $R$ \emph{partly constitutes} an entity $x$ if some component of $x$ consists, at least in part, in $R$'s holding, so that, were $R$ not to hold, that component would not exist, whatever its causal antecedents. $R$ merely \emph{causes} a component if $R$ figured in its production but its continued existence does not consist in $R$'s holding.
\end{definition}

\begin{definition}[Recognition]\label{d:rec}
A person $A$ \emph{recognizes} a person $B$ as having a status $S$ if $A$ treats $B$ as having $S$ in the relevant practices and holds herself and others accountable for doing so.
\end{definition}

\begin{definition}[Practical identity]\label{d:pi}
A \emph{practical identity} is a description under which a person values herself and finds her life worth living and her actions worth undertaking \citep{Korsgaard1996}, together with the standing that the description confers on her among others.
\end{definition}

\begin{assumption}[Components of the self]\label{a:comp}
The self, as characterized in PID-04-a, has, among other components, (i) a minimal subject, the first-person perspective of an organism on its own experience; (ii) the practical identities by which the person interprets and governs herself; and (iii) her standing as an interpreter of herself, that is, the authority of her self-interpretations.
\end{assumption}

\section{The argument}

\begin{lemma}\label{l:minimal}
The minimal subject is not constituted by recognition.
\end{lemma}

\begin{proof}
The first-person perspective of an organism on its experience consists in its experiences being had by it and available to it as its own. An organism whose experiences had this character would have it whether or not any other being treated it as anything. Recognition may be causally necessary for the development of a reflective grasp of oneself as a subject, as Mead held, but by Definition~\ref{d:const} that dependence is causal, not constitutive.
\end{proof}

\begin{lemma}\label{l:pi}
Practical identities are partly constituted by recognition.
\end{lemma}

\begin{proof}
A practical identity includes a standing among others (Definition~\ref{d:pi}). To be a parent, a physician, a citizen or someone's friend, in the sense that makes these descriptions available as reasons, is to hold a status, and holding such a status consists in part in being treated as holding it in the relevant practices: a physician whom no hospital, patient or licensing body treats as one cannot act as one, and a friend whose friend does not acknowledge the friendship is not one. MET-20-a shows that statuses of this kind exist in virtue of collective attitudes; the recognition of Definition~\ref{d:rec} is the form those attitudes take toward an individual. So, by Definition~\ref{d:const}, recognition partly constitutes practical identities.
\end{proof}

\begin{lemma}\label{l:standing}
A person's standing as an interpreter of herself is partly constituted by recognition.
\end{lemma}

\begin{proof}
A self-interpretation has authority, in the sense in which it governs the person's conduct and others' treatment of her, only to the extent that it receives uptake from others; this is the uptake principle by which PID-09-a shows narrative identity to be co-authored. To have standing as an interpreter of oneself is for one's self-interpretations to be presumptively accepted as such. That presumption is extended by others and can be withdrawn, as it is from those whose testimony about themselves is systematically discounted. So the standing consists in part in recognition.
\end{proof}

\begin{theorem}\label{t:main}
The self is partly, and only partly, constituted by relations of recognition.
\end{theorem}

\begin{proof}
By Assumption~\ref{a:comp}, practical identities and self-interpretive standing are components of the self. By Lemmas~\ref{l:pi} and~\ref{l:standing}, they are partly constituted by recognition, so the self is partly constituted by it. By Lemma~\ref{l:minimal}, the minimal subject, also a component, is not; so the self is not wholly constituted by recognition.
\end{proof}

Theorem~\ref{t:main} fits PID-04, on which the self is a real construction: recognition is one of the constructing activities. It fits PID-09, on which narrative identity is co-authored: Lemma~\ref{l:standing} identifies the relation through which the co-authors exercise their share.

\section{Objections and replies}

\begin{objection}[Merely causal]
Others shape us, but only as causes. Once formed, a self is what it is, whatever others now think.
\end{objection}

\begin{reply}
For the minimal subject, we agree (Lemma~\ref{l:minimal}). For practical identities, the counterfactual test of Definition~\ref{d:const} gives the opposite result: withdraw the recognition and the status ceases, however it was produced. A deposed judge retains her history and skills, but she is no longer a judge.
\end{reply}

\begin{objection}[Hostage to others]
If the self is constituted by recognition, others can unmake a person at will, which is both false and dangerous.
\end{objection}

\begin{reply}
The constitution is partial and the recognition is distributed. No single person can unmake another's practical identities, which are recognized in many practices, and a person can seek recognition in another community when one withholds it. That others can damage a self by misrecognition is not a reductio but a fact, the one \citet{Taylor1994} and \citet{Honneth1995} take as the starting point of a politics of recognition.
\end{reply}

\begin{objection}[The solitary]
A hermit who has lived alone for decades has a self, though no one now recognizes her.
\end{objection}

\begin{reply}
She has a minimal subject in full, and the practical identities and standing she once had persist as long as they are not withdrawn, in the way MET-20-a argues that statuses persist in the absence of collective rejection. What she lacks are identities she has never been recognized as having, and the view correctly predicts that, if she has come to interpret herself in ways no one has ever recognized, those interpretations have less standing than her others. The solitary has a self that is less social, not one that is not a self.
\end{reply}

\section{Conclusion}

The self is social in part. Its minimal core owes others at most its development; its practical identities and the authority of its self-interpretations consist in part in being recognized, and would not survive the withdrawal of recognition. Since both kinds of component belong to the self, the self is partly constituted by relations of recognition, and only partly.

\begin{thebibliography}{99}\small
\bibitem[Hegel(1807)]{Hegel1807} Hegel, G.~W.~F. (1807). \emph{Phenomenology of Spirit}. Cited by paragraph of the translation by A.~V. Miller (Oxford University Press, 1977).
\bibitem[Honneth(1995)]{Honneth1995} Honneth, A. (1995). \emph{The Struggle for Recognition: The Moral Grammar of Social Conflicts}, trans. J. Anderson. Polity Press.
\bibitem[Korsgaard(1996)]{Korsgaard1996} Korsgaard, C.~M. (1996). \emph{The Sources of Normativity}. Cambridge University Press.
\bibitem[Mead(1934)]{Mead1934} Mead, G.~H. (1934). \emph{Mind, Self, and Society}. University of Chicago Press.
\bibitem[Strawson(1997)]{GStrawson1997} Strawson, G. (1997). The self. \emph{Journal of Consciousness Studies}, 4(5--6), 405--428.
\bibitem[Taylor(1994)]{Taylor1994} Taylor, C. (1994). The politics of recognition. In A. Gutmann (ed.), \emph{Multiculturalism: Examining the Politics of Recognition} (pp.~25--73). Princeton University Press.
\end{thebibliography}

\end{document}
